% TRANSLATION UNIT: translate this file only.
% ===========================================================================
%  12 — Despliegue
%  The \section title is the canonical (Spanish) chapter name; the body prose
%  is written in ENGLISH and translated to Spanish later.
% ===========================================================================

% EVIDENCE INVENTORY (every claim below maps to a real artifact):
% EVIDENCE: four verbs, no manual step between them → dev/build:20-26 (builds, prints the next verb), dev/up:15-38 (raises the topology, waits on health checks), dev/migrate:1-95, dev/down:9-27.
% EVIDENCE: the composition forbids bypassing the verbs → dev/docker-compose.yaml:3-6.
% EVIDENCE: golang-migrate is the single migration tool, run as a container on the compose network → dev/migrate:4-8,36,66-71.
% EVIDENCE: configuration only by environment variables, discovery by hostname → dev/docker-compose.yaml:10-13.
% EVIDENCE: service counts by role → dev/docker-compose.yaml:186,204,234,255,276,314,338,351 (eight services), :387 (frontend), :13-102 (seven PostgreSQL), :108 (Redis), :120 (RabbitMQ), :137-180 (MailHog, Swagger UI, Adminer: development only); dev/up:15.
% EVIDENCE: e2e assertion budget doubled from 5 s to 10 s because one spec failed one run in three when the rest of the suite competed for the gateway → frontend/playwright.config.ts:29-36.
% EVIDENCE: the load run shares one host with seven databases, the broker, the cache and eight services → deploy/loadtest/README.md:198.
% EVIDENCE: dev compose builds frontend/Dockerfile.dev, whose entrypoint is `npm run start` = `ng serve` → dev/docker-compose.yaml:387; frontend/Dockerfile.dev:30; frontend/package.json:9.
% EVIDENCE: production image is a multi-stage Angular build served as static files by nginx-unprivileged, no Node in the final layer → frontend/Dockerfile:14-45,56,177-181.
% EVIDENCE: SPA fallback → frontend/Dockerfile:132; gzip on, level 6, types excluding woff2 → :75-95; no brotli because the image ships no ngx_brotli module → :63-71; dedicated `location = /robots.txt` → :110-121; cache headers → :100-109.
% EVIDENCE: main.js 77,328 B → 23,022 B gzip, CSS 21,588 B → 4,269 B gzip, Roboto 488,026 B TTF → 66,544 B woff2 → Lighthouse and served-payload measurement against frontend/Dockerfile, 2026-11-03; the JavaScript and font figures repeat chapter 11's load table.
% EVIDENCE: the config.js write permission failure existed only on the production path → frontend/Dockerfile:165-178.
% EVIDENCE: ci.yml jobs — contracts (buf lint + breaking) :19-31; Go matrix of five services, gofmt + golangci-lint + `go test -race`, audit adds testcontainers integration tests :33-76; Rust fmt/clippy/cargo test :78-98; TypeScript for learning, notification, frontend :100-130; dev workflow :135-160; documentation/script coherence :164-200.
% EVIDENCE: no job in ci.yml runs the prose barrier or the design-debt check (that absence is what this chapter found); the barriers now run in a separate workflow → .github/workflows/barriers.yml:1-60, added by commit 7112f09.
% EVIDENCE: barriers.yml triggers on every push on any branch and on every pull request, runs the prose barrier, the citation verifier and `design-debt.mjs --check`, and self-tests the barrier against a scratch file containing `behaviour` and `analysed` → .github/workflows/barriers.yml:12-20,30-57.
% EVIDENCE: only ERRORs fail the barrier; WARNs do not, and the citation verifier warns instead of failing for a work not read in full → .github/workflows/barriers.yml:31-39; docs/informe/scripts/prose-barrier.mjs:44; docs/informe/scripts/verify-citations.mjs:11-18.
% EVIDENCE: 39 of 46 sources read in full, so 7 are not → docs/informe/refs/fuentes-acceso.md:1-6.
% EVIDENCE: ci.yml still triggers on pushes to main and 001-* only, and on pull requests to main → .github/workflows/ci.yml:5-7.
% EVIDENCE: two VPS, 4 GB RAM / 2 vCPU each, application on the one with a public address and data on the one without → deploy/vps/README.md:3,14-24.
% EVIDENCE: one domain, Caddy routes /api/* to the gateway with the prefix stripped and everything else to the frontend → deploy/vps/Caddyfile:29-36; deploy/vps/README.md:35-38.
% EVIDENCE: two TLS issuers (ACME first, internal fallback) because the campus perimeter filtered the ACME validators → deploy/vps/Caddyfile:19-23; deploy/vps/README.md:604-636.
% EVIDENCE: 31 migrations across seven databases in 3.9 s on empty databases → deploy/vps/README.md:363-372.
% EVIDENCE: migrate down reverts the whole chain of the seven databases and deletes their data → deploy/vps/README.md:379-381; dev/migrate:26-33.
% EVIDENCE: no backups, one replica per service, Redis without authentication → deploy/vps/README.md:720-732.
% EVIDENCE: Kubernetes base plus dev/prod overlays, two replicas minimum, HPA to twenty, secrets referenced by name → deploy/k8s/README.md:3-13,20-26; deploy/k8s/base/users.yaml:12; deploy/k8s/overlays/prod/autoscaling.yaml:1-8.
% EVIDENCE: the production overlay does not fit the two VPSs before a single user arrives → deploy/vps/README.md:4-7.
% EVIDENCE: load profile 200 VUs, 1 min ramp, 3 min hold, 30 s down, one account per VU → deploy/loadtest/README.md:89-100.
% EVIDENCE: 64,268 requests over 4 m 08 s, 0.00 % HTTP failures, read p95 115 ms, action p95 763 ms → k6 summary of the attributed run, 2026-11-03 (chapter 11, load table).
% EVIDENCE: the edge limits per IP and per authenticated user at 600 requests per minute → services/api-gateway/cmd/gateway/main.go:204-208; deploy/loadtest/README.md:73-85.
% EVIDENCE: 160 responses were 429 within 58 ms, all before the measured window, and setup() now uses per-account synthetic addresses in 10.255.0.0/16 → deploy/loadtest/README.md:160-186; deploy/loadtest/k6-scenarios.js:568-574.

\section{Despliegue}

% DONE: deployment as the fourth artifact of the system — the local composition, the
% image that ships, the server configuration, the pipeline, and the operations that
% exist against those that only exist as files. Every figure is attributed.

\subsection{Local reproducibility as a design goal}

The first deployment decision in FintCart was not where the system runs but whether a second person
can raise it. The repository answers with four verbs: \code{dev/build} builds the development images,
\code{dev/up} raises the topology and blocks until it reports healthy, \code{dev/migrate} applies every
service's migrations, and \code{dev/down} stops the composition while keeping the volumes unless asked
to delete them. No step between them is manual, and the composition file states the rule in its header: operate always
with the verbs, and if something cannot be done with the scripts, that is a defect of the scripts.

Five stacks live in the monorepo, and each stack's own tool — Go migrations for Go, an ORM runner for
NestJS, a SQLx macro for Rust — would mean five ordering semantics and five ways to be half-applied. FintCart uses one tool, \code{golang-migrate}, for all five stacks, executed as a
container attached to the compose network so that no machine needs it installed. The consequence is
that \enquote{the schema is at version N} is one statement across seven databases.

Configuration follows the same line: everything arrives by environment variable, and every address is
a hostname resolved inside the compose network. The frontend is the interesting case: a compiled SPA
cannot read the environment, because its JavaScript runs in the browser. Baking the API URL into the
bundle is the usual answer, and it creates one image per environment — the artifact tested in staging
is then not byte-for-byte the artifact served in production. FintCart writes \code{config.js} at
container start from \code{API_BASE_URL}, so one image serves every environment and the difference
between them is a variable.

Why insist? Because \enquote{it works on my machine} is not a convenience problem but an undocumented
part of the specification. A system only its author can raise has an implicit dependency graph — a
locale, a port already taken, a tool at the right version — that no artifact records, and the failure
surfaces later as an environment difference rather than as the missing declaration it is. An
unreproducible deployment also cannot be measured: every number in Chapter 11 describes a system
raised the same way each time, and if the raising varied, the numbers would describe machines.

The claim is worth only what enforces it. A dedicated CI job checks out the repository, runs
\code{dev/build}, \code{dev/up}, and \code{dev/migrate up}, then reverts the migrations and reapplies
them, then runs \code{dev/down}; if any verb fails, the pipeline fails. That job is the difference
between a promise and a gate.

\subsection{Container topology and what it costs}

Table~\ref{tab:topology} lists what \code{dev/up} raises: eight backend services, the SPA, seven
PostgreSQL instances, one Redis, one RabbitMQ broker, and three tools that exist only in development.

\begin{table}[htbp]
  \centering
  \begin{tabular}{llp{7.4cm}}
    \hline
    Component & Count & Role in the composition \\
    \hline
    Backend services & 8 & Gateway (REST edge, no state), auth, users, learning, simulator, orchestrator, notification, audit \\
    SPA & 1 & Angular bundle; development server locally, nginx in production \\
    PostgreSQL 16 & 7 & One database per service with state; the gateway has none \\
    Redis 7.4 & 1 & JWT blacklist and refresh tokens (auth), rate limiting (gateway) \\
    RabbitMQ 4.0 & 1 & Producer to consumer events; the gateway never touches it \\
    Development tools & 3 & MailHog, Swagger UI, Adminer: no production counterpart \\
    \hline
  \end{tabular}
  \caption{Topology raised by \code{dev/up}.}
  \label{tab:topology}
\end{table}

Seven databases for eight services is not an oversight: the gateway owns no domain state, so the count
follows the services that own data. The table restates the constitutional rules in operational form: a
violation would show up here before it reached an incident.

The cost belongs here too. The full composition puts seven PostgreSQL instances, a broker, a cache,
eight services, a SPA, and the test harness in competition for one host's cores. That contention was
measured, not assumed: one end-to-end spec failed one run in three when the rest of the suite competed
for the gateway and passed first time in isolation, so the assertion budget went from five to ten
seconds with the reason recorded rather than absorbed silently. A multi-second stall in a browser
assertion is the shape this contention takes. It does not invalidate the composition — the composition
exists to be raised, not to be fast — but the host is a test rig, and a rig that stalls changes what a
test measures.

\subsection{The image that ships is not the image that develops}

The development composition builds \code{frontend/Dockerfile.dev}: it runs \code{npm run start},
which is \code{ng serve}, a server that recompiles on change, keeps source maps, and serves
unoptimized output. The production file compiles the bundle in a Node stage, copies
the static output into \code{nginx-unprivileged}, and contains no Node at all.

This distinction produced a false measurement. An early performance reading was taken against the
stack raised by \code{dev/up} — that is, against \code{ng serve} — and the numbers looked like a slow
application. They described a development server: uncompressed modules, no minification, a watcher
beside the process being measured. None of it could have been true of what a student downloads. The
measurements in Chapter 11 were therefore taken against \code{frontend/Dockerfile}, with the served
bundle's hash compared against the build output so the payload measured was provably the code under
discussion. Measuring the artifact you ship sounds like a tautology; it has to be written
down because the convenient artifact is always the one already running.

The same split hid a second defect in the other direction. The production image writes
\code{config.js} at startup into a directory owned by the nginx user; before that directory existed,
the first real deployment failed with a permission error because the served tree belonged to root. No
number of runs in development could have caught it: that script never executes there. A development
path that skips a production code path is not a smaller deployment, it is a different program, and the
difference sits exactly on the unexercised lines.

\subsection{The web server is part of the deliverable}

What ships is static files and the configuration of the server that hands them out; four decisions
are load-bearing.

The first is \code{try_files $uri $uri/ /index.html}. Without it, reloading the browser on a
client-side route returns 404, because the route exists in the Angular router, not as a file. The
second is compression: the bundle was served without \code{Content-Encoding} although the browser
requested it, so plain text crossed the network at full size. gzip is enabled at level six, and the
types list excludes woff2, png, and jpg, already compressed. Brotli is absent, and the absence is
verified against the image's compiled module list: the \code{nginx-unprivileged:alpine} image ships no
\code{ngx_brotli} module, so a \code{brotli on;} directive is rejected at startup and the container
never comes up.

The third is a dedicated \code{location = /robots.txt}: with only the SPA fallback, a request for that
file received \code{index.html} with status 200 and an HTML content type, which an audit reports as an
invalid robots file. The explicit location serves the real file and returns an honest 404 if it ever
leaves the build. The fourth is caching: hashed names are immutable for a year, while
\code{index.html} is \code{no-cache} because it is the file that points at the new hashes.

Table~\ref{tab:payload} reports what the first two decisions are worth.

\begin{table}[htbp]
  \centering
  \begin{tabular}{llrr}
    \hline
    Asset & Mechanism & Before & Served \\
    \hline
    \code{main.js} & gzip, level 6 & 77,328 B & 23,022 B \\
    CSS bundle & gzip, level 6 & 21,588 B & 4,269 B \\
    Roboto & woff2 instead of TTF & 488,026 B & 66,544 B \\
    \hline
  \end{tabular}
  \caption{Static payload before and after the server and format fixes.}
  \label{tab:payload}
\end{table}

The JavaScript figure is a seventy per cent reduction and the CSS figure eighty per cent, both
produced by a server directive and neither by a line of application code. The font row is a different
mechanism and must not be added to the same total: 488,026 bytes of TrueType became 66,544 bytes of
woff2 because the build emits a compressed format, not because nginx compressed the old one. The last
seventy per cent of the transfer was therefore decided in the server configuration, which a project
that stops at the bundle leaves to whoever configures the server — who, in most projects, does not
exist.

\subsection{Continuous integration}

The pipeline runs six things, and naming them matters, because \enquote{CI runs the tests} would
hide both what is covered and what is not.

Contracts come first: \code{buf} lints the proto tree and checks backward compatibility against
\code{main}, so an incompatible message change fails before any service compiles. Then Go, as a
matrix over the five Go services: format check, \code{golangci-lint} with the repository
configuration that also bans float for money, and \code{go test -race}; audit additionally runs its
testcontainers integration tests against a real PostgreSQL. Then Rust: \code{cargo fmt --check},
\code{cargo clippy --all-targets -D warnings}, and \code{cargo test --all-targets}. Then TypeScript,
as a matrix over learning, notification, and the frontend: \code{npm ci}, lint, build, and test. Then
the development workflow job described earlier. Then a coherence job that fails if the documentation
cites a \code{dev/} verb that does not exist or omits one that does.

The prose barrier, the citation verifier, and the
design-debt script existed and exited non-zero on failure, but no job in \code{ci.yml} ran them: they
ran when someone remembered. A check that depends on memory is weaker than one that fails the build;
its failure mode is silence, and a rule nobody hears broken stops being a rule.

Writing this chapter found that gap, and commit \code{7112f09} closes it. A separate workflow,
\code{.github/workflows/barriers.yml}, runs the three scripts on every push, on any branch, and on
every pull request — separate from \code{ci.yml} on purpose, because they are Node and files and
should not drag \code{cargo test} into a documentation push. It also tests the tester: against a
scratch file with \code{behaviour} and \code{analysed}, the barrier must exit non-zero, since a
verifier that has never failed has not shown it can fail. Severity stays deliberately uneven: only
ERRORs stop the build, so the citation verifier warns rather than fails when a cited work was not read
in full — seven of the forty-six sources are not — and the prose barrier's warnings do not fail
either, because a declared debt is not a build error.

\subsection{Operations that exist, and operations that do not}

What exists is a deployment on two virtual machines of four gigabytes and two cores each. The
application machine, the only one with a public address, runs the eight services, the frontend, and
Caddy; the data machine runs the seven databases, Redis, and RabbitMQ, reachable only over the private
network. Caddy terminates TLS on a single domain and routes by path — \code{/api/*} to the gateway
with the prefix stripped, everything else to the frontend — which lets one certificate cover both.
TLS has two issuers, ACME and Caddy's internal authority as fallback, because the campus perimeter
filtered the ACME validators: with ACME alone the site had no certificate, and with the authority
alone someone had to remember to undo the workaround.

Migrations ran as a separate step before the new code started: thirty-one across seven databases in
3.9 seconds. That figure is a floor: the databases were empty, and migrations that walk whole tables grow with the
rows. Rollback deserves the same honesty: \code{dev/migrate
down} does not revert the last migration, it reverts the entire chain of the seven databases and
deletes their data, so real rollback begins with a dump taken before migrating, and schema and code
move back together.

What does not exist is the list that usually fills this section. There is no Kubernetes in production.
The repository holds a kustomize base and two overlays with deployments, probes, autoscalers, and
disruption budgets, and the production overlay asks for two replicas per service and up to twenty
under load; on two machines of four gigabytes those requests do not fit before a single student
arrives, which is why the compose tree exists. We did not run the manifests against a cluster, so this
chapter does not present them as a running system: their content can be inspected, their behavior
cannot be certified here. There is no autoscaling, no blue-green switch, no
second environment; one replica per service means a restart is a visible outage; the seven volumes
have no automated backup; Redis has no password, protected only by the network boundary.

One deployment was operated and several were described; enumerating promises with no run behind them
would be a catalog.

\subsection{Load testing as a deployment concern}

The load profile is a property of the deployment, not of the code. Two hundred virtual users was not
chosen for realism: it was the size of the account pool, and one account per virtual user was forced
by the edge, which limits per IP and per authenticated identity at 600 requests per minute each. The
reported run held that profile for three minutes and issued 64,268 requests over four minutes and
eight seconds with no HTTP failures, a read p95 of 115 ms, and an action p95 of 763 ms.

Those numbers are honest only with their ceiling attached. The run happened on a single development
host sharing its cores with the whole composition, so it is not a production capacity figure. The
ceiling itself belongs to the limiter: 600 requests per minute per client is ten per second per
client, a bound the aggregate of two hundred independent clients never approaches. Its setup traffic
left from one address, and the access log
attributes 160 rejections inside 58 milliseconds to that phase, all of them before the measured window
opened; the fix was a synthetic source address per account, and the breakdown by phase is now recorded
so a future run can say whether a failure belongs to the platform or to the harness.

What the run supports is narrow: two hundred clients, each inside its allowance, served for three
minutes without an error and with percentiles far below the thresholds. It supports no
statement about how many students the platform holds, because multiplying the observed rate would be
arithmetic on a number the limiter fixed. A capacity figure needs the replicas the manifests describe
and hardware the two machines do not have.
